% ============================================================
% Prefacio — traducción del original "Préface" (1889)
% ============================================================
\chapter*{Prefacio}
\addcontentsline{toc}{chapter}{Prefacio}
\markboth{Prefacio}{Prefacio}

La Óptica es la parte más avanzada de la física; la teoría llamada de las
ondulaciones forma un conjunto verdaderamente satisfactorio para el
espíritu; pero no hay que pedirle lo que no puede darnos.

Las teorías matemáticas no tienen por objeto revelarnos la verdadera
naturaleza de las cosas; sería esa una pretensión irrazonable. Su único
fin es coordinar las leyes físicas que la experiencia nos da a conocer,
pero que sin el auxilio de las matemáticas ni siquiera podríamos enunciar.

Poco nos importa que el éter exista realmente; eso es asunto de los
metafísicos; lo esencial para nosotros es que todo ocurra como si
existiera, y que esta hipótesis sea cómoda para la explicación de los
fenómenos. Después de todo, ¿tenemos acaso otra razón para creer en la
existencia de los objetos materiales? Tampoco esa es más que una hipótesis
cómoda; solo que esta nunca dejará de serlo, mientras que sin duda llegará
un día en que el éter sea rechazado por inútil.

Pero ese mismo día, las leyes de la óptica y las ecuaciones que las
traducen analíticamente seguirán siendo verdaderas, al menos como primera
aproximación. Será, pues, siempre útil estudiar una doctrina que enlaza
entre sí todas esas ecuaciones.

Las teorías propuestas para explicar los fenómenos ópticos mediante las
vibraciones de un medio elástico son muy numerosas e igualmente
verosímiles. Sería peligroso limitarse a una sola de ellas; se correría
así el riesgo de depositar en ella una confianza ciega y, por tanto,
engañosa. Hay que estudiarlas todas, y es precisamente la comparación lo
que puede resultar más instructivo.

Desgraciadamente, para conocerlas todas es necesario recurrir a las
memorias originales, a menudo difíciles de encontrar y penosas de leer. Al
pasar de una a otra, todo cambia: la notación, las maneras de pensar y de
expresarse, y la comparación se vuelve casi imposible.

He creído, pues, que no sería inútil reunir en un pequeño volumen, y en un
cuadro de conjunto, estas diversas doctrinas, publicando las lecciones que
he profesado en la Sorbona en 1887-1888. Estoy muy agradecido al señor
Blondin, quien hizo posible mi tarea recogiendo y redactando estas
lecciones.

Para leer este libro con provecho, es necesario poseer, al menos en sus
rasgos generales, las leyes experimentales de la óptica física. En
efecto, no habría podido recorrer en un solo semestre una extensión tan
grande de materias si no hubiera sabido que estas leyes ya eran familiares
a mis oyentes, licenciados en su mayoría en ciencias físicas.

La teoría de las ondulaciones descansa sobre una hipótesis molecular; para
unos, que creen descubrir así la causa bajo la ley, esto es una ventaja;
para otros, es un motivo de desconfianza; pero esta desconfianza me parece
tan poco justificada como la ilusión de los primeros.

Estas hipótesis no desempeñan sino un papel secundario. Habría podido
sacrificarlas; no lo hice porque la exposición habría perdido en claridad,
y esa fue la única razón que me lo impidió.

En efecto, de las hipótesis moleculares no tomo prestado más que dos
cosas: el principio de la conservación de la energía y la forma lineal de
las ecuaciones, que es la ley general de los pequeños movimientos, como de
todas las pequeñas variaciones.

Esto explica por qué la mayoría de las conclusiones de Fresnel subsisten
sin cambio cuando se adopta la teoría electromagnética de la luz. No
hablaré en este volumen de esta teoría, pues me reservo exponerla en
detalle en otra obra donde publicaré mis lecciones del segundo semestre.
He creído deber dedicarme primero a profundizar en las ideas de Fresnel;
esto me pareció la mejor preparación para el estudio del pensamiento de
Clerk Maxwell.

Una última palabra. En el capítulo relativo a la difracción, he
desarrollado ideas que creía nuevas. No mencioné a Kirchhoff, cuyo nombre
habría debido citarse en cada línea. Aún hay tiempo de reparar este olvido
involuntario; me apresuro a hacerlo remitiendo a los \emph{Sitzungsberichte}
de la Academia de Berlín (1882, segundo semestre, página 641).

\bigskip
\noindent\textit{París, 2 de diciembre de 1888.}

\bigskip
\noindent\textsc{H. Poincaré.}
